\begin{claim}{C.54}{Proved}{Two metabelian quotient images give exponential closure contractions}
Use the actual global quotient family of C.53 for the
ambient-null based plane loop $f_t$.
Let $O$ be an open $G$-orbit admitting two distinct
partitions in that family whose actual intrinsic target
images are metabelian. On the universal covering
$p:X\to O$, the group $N$ generated by the based source
lifts is closed in $\operatorname{Homeo}(X)$.
With that actual topology, it is one of the following
simply connected groups:
\[
\begin{gathered}
 \R^2,\qquad \operatorname{Aff}^+(\R),\qquad
 N_{a,b}=\R^2\rtimes\R\quad((a,b)\ne(0,0)),\\
 \operatorname{Aff}^+(\R)\times\operatorname{Aff}^+(\R).
\end{gathered}                                           \tag{C.54.1}
\]
In $N_{a,b}$ the real parameter $t$ acts on the two
translation coordinates by
$\operatorname{diag}(e^{at},e^{bt})$.
Thus zero, equal, unequal, and opposite nonzero
weights are retained.
Each group has the explicit globally invertible
exponential map given below.

The based source exponential fractions descend and
extend to a jointly continuous based contraction
through homeomorphisms of the entire $\overline O$,
with a jointly continuous inverse.
They satisfy every actual global quotient equation,
are independent of the selected pair of metabelian
quotients, and are transported exactly by $H$.
They agree on common orbit frontiers and with the
prescribed contraction of C.53 on $X^\dagger$.
Consequently they paste over $X^\dagger$ together
with any finite collection of these closures.
No compactness of $O/Q_O$ is required.
Uniform continuity over an arbitrary infinite
collection of the new closures is not asserted here.
\end{claim}

\begin{proof}
\emph{The actual target groups and the two coordinate domains.}
Choose two distinct real lifted quotient maps
$F_1,F_2:X\to\R$, with their original target covers
retained. Their based target paths $b_{i,t}$ are
loops, by C.29 for circle targets and directly
for line targets. As in C.52.4, the generated
groups $T_i=\langle b_{i,t}\rangle$ are transitive
on the real lines. The exact real-target
comparison in C.52 makes each $T_i$ a full
translation group or a full positive affine
group after a line conjugacy. These are the
actual groups with their intrinsic
homeomorphism-group topologies.

C.45 identifies the pair of quotient coordinates
with one of the three domains
$\R^2$, $\{x<y\}$, or $\{x<y<x+1\}$, by the
specified independent line reparametrizations.
The based source path has the coordinate form
$(b_{1,t},b_{2,t})$ and is a closed loop, since
the target endpoints are the identity and the
pair map is injective.

In the half-plane domain the two coordinate
paths are identical, and their generated group
is the same $T$. A transitive abelian line
group is conjugate to translations. Its
diagonal action cannot be transitive on an
open two-dimensional domain: in translation
coordinates the difference of the two
coordinates is constant along each orbit.
Thus $T$ is the full affine group.
Conjugate both coordinates by the same line
homeomorphism. The source action becomes
\[
          (x,y)\longmapsto(e^t x+u,e^t y+u),
                         \qquad x<y.                  \tag{C.54.2}
\]
The exact source change
$(x,y)\mapsto(x,\log(y-x))$ identifies this with
the simply transitive plane action
\[
                    (x,r)\longmapsto(e^t x+u,r+t).
                                                               \tag{C.54.3}
\]
Its actual generated group is the full
$\operatorname{Aff}^+(\R)$.

The band domain is impossible under the
present metabelian-target hypothesis.
Here again the two paths and their group $T$
coincide, and $T$ commutes with unit translation.
The full affine line group has trivial
centralizer: a map commuting with every
translation is a translation, and commuting
with a nontrivial dilation makes that
translation the identity.
Thus $T$ cannot be affine.
If $T$ is conjugate to translations, the
commuting unit translation is conjugate to
translation by a fixed nonzero constant.
The band therefore becomes a constant-width
band in those same coordinates, and the
diagonal translation action again preserves
the coordinate difference. It is not
transitive on the band, a contradiction.

It remains to classify the full-plane case.
Conjugate the two target groups independently
to their standard translation or affine
actions. The actual source group is then a
path-connected subgroup of
\[
 A=\operatorname{Aff}^+(\R)^2
   =\{(s_1,s_2,u_1,u_2):
        (x_1,x_2)\mapsto
             (e^{s_1}x_1+u_1,e^{s_2}x_2+u_2)\}.
                                                               \tag{C.54.4}
\]
We do not assume that this subgroup is closed.

The based source group is transitive on $X$.
Indeed its compact path-and-inverse set $P$
and the countable deck group $\Delta$ satisfy
$X=\bigcup_{\delta\in\Delta,n\geq1}\delta P^n x$,
as in C.49.3. Baire gives an open orbit.
Every other orbit is a deck translate,
because the based source path commutes with
every source deck map. Hence all orbits are
open, and connectedness of the covering plane
gives transitivity. Yamabe's theorem now
identifies this path-connected subgroup with
the analytic subgroup of its Lie algebra
\cite[pp.~13--14]{C_Yamabe1950}. We establish
closedness and the induced topology separately.

\emph{All transitive analytic subgroups of the product affine action.}
Write the Lie algebra of (C.54.4) as
$\mathfrak a=D\ltimes V$, where $D=\R^2$ is
the pair of log-scale directions, $V=\R^2$
is the translation plane, and
\[
 [(d,v),(e,w)]
       =(0,(d_1w_1-e_1v_1,\ d_2w_2-e_2v_2)).
                                                               \tag{C.54.5}
\]
Let $\mathfrak n$ be the generated group's
Lie algebra and put
\[
 S=\operatorname{pr}_D(\mathfrak n),\qquad
 W=\mathfrak n\cap V.
                                                               \tag{C.54.6}
\]
The subspace $W$ is invariant under the
diagonal action of $S$.
Transitivity makes the tangent evaluation
of $\mathfrak n$ onto the source plane
surjective at every point.

For $\dim S=0$, this forces $W=V$ and gives
all translations.
For $\dim S=1$, choose a lift
$Z=((a,b),v_0)$ of a nonzero generator of $S$.
The case $W=0$ has dimension one and cannot
be transitive. If $W=V$, subtracting $v_0$
from $Z$ gives exactly the Lie algebra and
connected group $N_{a,b}$ in (C.54.1).

Suppose instead $\dim W=1$. If $a=b\ne0$,
the induced affine vector field on $V/W$
has nonzero scalar linear part and hence
a fixed point. The corresponding affine
line in $V$ is invariant under the entire
connected subgroup, contradicting
transitivity.
If $a\ne b$, diagonal invariance makes $W$
one of the two coordinate axes.
For $W=\R(1,0)$, the second-coordinate
vector field of $Z$ is $by+(v_0)_2$.
If $b\ne0$, its fixed level is an invariant
affine line, again impossible.
If $b=0$ and $(v_0)_2=0$, the second
coordinate is constant on every orbit.
The remaining possibility is
\[
                  b=0,\qquad a\ne0,\qquad (v_0)_2\ne0.
                                                               \tag{C.54.7}
\]
A translation of the first source coordinate
removes $(v_0)_1$, and replacing the second
coordinate by $ay/(v_0)_2$ gives exactly
(C.54.3). The case $W=\R(0,1)$ is its
coordinate exchange. These are all
one-dimensional $S$ cases.

For $\dim S=2$, invariance under all diagonal
matrices makes $W$ either zero, one coordinate
axis, or $V$. If $W\ne V$, choose a coordinate
$i$ on which every vector of $W$ is zero.
For a linear section of
$\mathfrak n\to S=D$, let $c_i(d)$ be its
translation coordinate there. This is
independent of the chosen section modulo $W$.
Equation (C.54.5), whose bracket must belong
to $W$, gives
\[
                  d_i c_i(e)-e_i c_i(d)=0
                                  \quad(d,e\in D).
                                                               \tag{C.54.8}
\]
Choosing $e_i=1$ shows
$c_i(d)=c\,d_i$ for one constant $c$.
Every Lie-algebra vector field then preserves
the level $x_i=-c$. Its flows, and hence the
connected analytic subgroup that they
generate, preserve that affine line.
This contradicts transitivity.
Thus $W=V$, and $\mathfrak n=\mathfrak a$,
giving the full product affine group.

\emph{Closed groups, actual topology, and logarithms.}
The preceding list consists of explicit closed
subgroups of $A$. For $N_{a,b}$, the scale
vector lies on the closed line
$\R(a,b)$ and both translations are arbitrary.
The simply transitive affine case (C.54.3)
has, in its displayed source coordinates,
the closed equations $s_2=0$ and $u_2=s_1$.
The translation plane and $A$ itself are
closed as well.
Each listed group is connected, with exactly
the Lie algebra just obtained. Uniqueness
of the connected analytic subgroup with a
specified Lie algebra therefore identifies
the original generated subgroup with that
explicit closed group.

The topology of $A$ in $\operatorname{Homeo}(\R^2)$
is its ordinary parameter topology.
Evaluation at $(0,0),(1,0),(0,1)$ recovers
the translations and the two positive scale
factors continuously. A limit which is
still a homeomorphism cannot have either
scale zero. Conversely convergent parameters
give compact-open convergence of the maps
and their inverses by the affine formulas.
Thus $A$ is closed in this homeomorphism
group and has the asserted actual topology.
The source conjugacies used above preserve
compact-open group topology. This proves
the same assertions for the actual $N$ on
$X$, rather than only for its analytic
Lie-group structure.

Here are the promised exponential maps.
Set
\[
 \sigma(z)=\int_0^1 e^{rz}\,dr
   =
   \begin{cases}(e^z-1)/z,&z\ne0,\\1,&z=0.\end{cases}
                                                               \tag{C.54.9}
\]
This function is positive for every real
$z$. In $N_{a,b}$ the exact group law is
\[
 (t,u,v)(t',u',v')
     =(t+t',u+e^{at}u',v+e^{bt}v'),
                                                               \tag{C.54.10}
\]
and
\[
 \begin{split}
 \exp(t,u,v)&=(t,\sigma(at)u,\sigma(bt)v),\\
 \log(t,u,v)&=(t,u/\sigma(at),v/\sigma(bt)).
 \end{split}                                             \tag{C.54.11}
\]
The positive affine group has the same
formula with one translation coordinate
and weight one. For the full product $A$,
\[
 \begin{split}
 \exp(s_1,s_2,u_1,u_2)
      &=(s_1,s_2,\sigma(s_1)u_1,\sigma(s_2)u_2),\\
 \log(s_1,s_2,u_1,u_2)
      &=(s_1,s_2,u_1/\sigma(s_1),u_2/\sigma(s_2)).
 \end{split}                                             \tag{C.54.12}
\]
The translation exponential is the identity.
These formulas follow by integrating the
displayed affine vector fields, or directly
by differentiating their one-parameter
subgroups. Positivity of $\sigma$ proves
that each is a global continuous inverse
pair, including every zero weight.
The displayed group parameter spaces are
Euclidean spaces, so every group in
(C.54.1) is simply connected.

\emph{The exact action on the closure and on all global quotients.}
Use the original global closure $L$ and
the specified source covering to form
\[
 E_O=\{(l,n)\in L\times N:pn=(l|_O)p\}.
                                                               \tag{C.54.13}
\]
This is a closed Polish group.
Every $n\in N$ is an actual finite word
in the based source lifts; the identical
word in the original plane loop gives a
preimage in $E_O$. Projection to $N$ is
therefore onto and open.
Its kernel fixes $O$ pointwise and, by
continuity, fixes $\overline O$ pointwise.
Consequently the continuous action
$(l,n,z)\mapsto lz$ descends through the
open product quotient
$E_O\times\overline O\to N\times\overline O$.
This is the actual jointly continuous
$N$-action on the closure. Every value
has a global word preimage, so its orbits
are precisely the original $G$-orbits
there, with the original labels retained.

Define the source family by
\[
          n_{s,t}=\exp_N\bigl((1-s)\log_N\widetilde f_t\bigr).
                                                               \tag{C.54.14}
\]
The source deck group commutes with every
generator, hence with $N$ itself.
Formula (C.54.14) descends through $p$
and acts on $\overline O$ through the
just-proved quotient action. It and its
inverse are jointly continuous and satisfy
all based-edge identities.

For any actual global quotient
$\Phi_z:O\to B$, equivariance makes the
kernel of (C.54.13) fix its entire image:
$l\Phi_z(w)=\Phi_z(lw)=\Phi_z(w)$.
The intrinsic restriction continuity
proved in C.42 therefore descends through
the same open quotient to a continuous
$N$-action on $B$ with image exactly
$G|_B$, as in C.53.7.
All groups in (C.54.1) are metabelian.
Thus this target image is one of the
actual Lie groups in C.37.9.
The continuous homomorphism, and its
unique based covering lift for a circle,
intertwine (C.54.14) with the exponential
branch of C.52.1. Hence every actual
global quotient equation holds, not only
the two equations selected at the start.

\emph{Boundary comparison, choices, and finite pasting.}
For a nontrivial nonopen $G$-orbit
$B\subset\overline O$, the continuous
$N$-action just constructed has exactly
that orbit. For $z\in B$, its orbit map
gives a continuous bijection
$N/N_z\to B$ with the ambient topology
on the image. Its manifold dimension is
one: dimension zero would give a point
because $N$ is connected, while dimension
at least two would give an injective
map of a small two-dimensional disk into
the plane with interior image, by
invariance of domain. This would make
$B$ contain an ambient open subset,
contrary to its being nonopen.

On each compact interval of this
one-manifold the orbit map is an arc.
C.28 identifies its intrinsic interval
with that same compact arc. Consequently
the orbit bijection is a homeomorphism
onto the intrinsic line or circle $B$.
The induced action is continuous in
that intrinsic topology, and its
compact-open group topology, by the
continuous-action exponential law.
Its image is exactly $G|_B$ and is
metabelian. C.37.9 and the same based
covering comparison therefore identify
the boundary value of (C.54.14) with
the precise exponential branch of
C.52.1. Fixed boundary points remain
fixed.

The group $N$ was defined by the actual
based source lifts, independently of
the chosen quotient pair. The proof
has identified its actual topology and
its unique global exponential and
logarithm. A different quotient pair
therefore gives the identical family.
Changing the source covering conjugates
the based paths, with their specified
identity initial lift, and hence
conjugates the identical formula.
An operator $h\in H$ gives this same
comparison between the source covers
of $O$ and $hO$. Its lifted conjugacy
identifies the actual generated groups
and their paths. Lie homomorphism
naturality and uniqueness of logarithms
give exact transport of (C.54.14).

Distinct open-orbit closures meet only
in nonopen orbits. Their intrinsic
images and original based target
paths are the same $G|_B$, so the
exponential values just proved agree
there, including the full real
circle lifts.
On an open orbit already in
$U^\dagger$, the earlier construction
uses this same actual group generated
by the based source lifts.
Thus the two formulas agree on that
open orbit and its closure. For a
new orbit disjoint from $X^\dagger$,
its intersection with $X^\dagger$
consists only of nonopen orbits,
where both values are again the
same exponential prescription.
Finite closed-set pasting applies
to the maps and their joint inverses.
This proves the stated finite relative
contraction without a cocompactness
assumption. No uniform word bound for
the new three- or four-dimensional
groups has been used or asserted.
\end{proof}
